\chapter{微分方程}\label{ch:g12:diffeq}

微分方程刻画函数与其导数之间的关系。物理学、
化学、生物学和经济学都以这种形式表述其规律：物体冷却的速率、
放射性原子核的衰变、种群的增长，都是关于 $y'$ 的陈述。
本章完整求解
带常系数的一阶线性微分方程。

\section{方程 \texorpdfstring{$y' = ay$}{y' = ay}}

\begin{definition}[微分方程]\label{def:g12:diffeq:ode}
称 \emph{微分方程}\index{微分方程} 为一个方程，
其未知量为函数 $y$，涉及 $y$ 及其导数。
区间 $I$ 上的\emph{解}是指一个可导函数，
它在 $I$ 的每一点都满足该方程。
\end{definition}

\begin{theorem}[$y' = ay$ 的解]\label{thm:g12:diffeq:homogeneous}
设 $a \in \R$。在 $\R$ 上，方程 $y' = ay$ 的解
恰为下列函数：
\[
y(x) = C\,\eu^{ax}, \qquad C \in \R .
\]
对每一对 $(x_0, y_0)$，存在唯一的解
满足 $y(x_0) = y_0$。
\end{theorem}

\begin{proof}
每个 $y = C\eu^{ax}$ 都满足 $y' = Ca\,\eu^{ax} = ay$。反之，设 $y$
为任一解，并令 $z(x) = y(x)\,\eu^{-ax}$。则
\[
z'(x) = y'(x)\,\eu^{-ax} - a\,y(x)\,\eu^{-ax}
= \bigl(y'(x) - a y(x)\bigr)\eu^{-ax} = 0,
\]
于是 $z$ 为常数，记为 $C$，从而 $y(x) = C\eu^{ax}$。初始条件
$y(x_0) = y_0$ 唯一确定 $C = y_0 \eu^{-a x_0}$。
\end{proof}

\begin{remark}
指数函数因此由所有微分方程中最简单的一个\emph{刻画}：
增长量与大小成正比。这正是它
在自然界中无处不在的原因。
\end{remark}

\begin{omfigure}
\begin{tikzpicture}[scale=1.05]
  \foreach \x in {-2.8,-2.4,...,2.9}
    \foreach \y in {-1.8,-1.4,...,1.9}
      \draw[gray!60, thin] (\x,\y) ++({atan(\y)}:0.11)
        -- ++({atan(\y)+180}:0.22);
  \draw[-Stealth] (-3.2,0) -- (3.4,0) node[below, font=\small] {$x$};
  \draw[-Stealth] (0,-2.2) -- (0,2.4) node[left, font=\small] {$y$};
  \draw[omDef, thick] plot[domain=-3:0.69, samples=60] (\x, {exp(\x)});
  \draw[omThm, thick] plot[domain=-3:1.89, samples=60] (\x, {0.3*exp(\x)});
  \draw[omProp, thick] plot[domain=-3:1.38, samples=60] (\x, {-0.5*exp(\x)});
\end{tikzpicture}

{\small 方程 $y' = y$ 在平面每一点规定了一个斜率
（灰色线段）。解 $C\eu^{x}$ --- 这里 $C = 1$（蓝色）、
$C = 0.3$（红色）、$C = -0.5$（橙色） --- 恰为沿此斜率场
走向的曲线。}
\end{omfigure}

\begin{example}[放射性衰变]\label{ex:g12:diffeq:decay}
某放射性量满足 $N' = -\lambda N$，其中 $\lambda > 0$，于是
$N(t) = N_0\,\eu^{-\lambda t}$；关于半衰期的讨论见
\cref{ex:g12:exp:halflife}。
\end{example}

\section{方程 \texorpdfstring{$y' = ay + b$}{y' = ay + b}}

\begin{theorem}[$y' = ay + b$ 的解]\label{thm:g12:diffeq:affine}
设 $a \neq 0$，$b \in \R$。在 $\R$ 上，$y' = ay + b$ 的解
恰为下列函数：
\[
y(x) = C\,\eu^{ax} - \frac{b}{a}, \qquad C \in \R .
\]
常值函数 $y_p = -\frac ba$ 称为 \emph{平衡解}。
对每一 $(x_0, y_0)$，存在唯一的解满足 $y(x_0) = y_0$；若
$a < 0$，则每个解都趋向平衡解 $-\frac ba$，
当 $x \to +\infty$ 时。
\end{theorem}

\begin{proof}
常数 $y_p = -\frac{b}{a}$ 满足 $y_p' = 0 = a y_p + b$。现设 $y$
为解的充要条件是
\[
(y - y_p)' = y' = ay + b = a(y - y_p) + \underbrace{a y_p + b}_{=\,0}
= a (y - y_p),
\]
即 $z = y - y_p$ 满足 $z' = az$。
由\cref{thm:g12:diffeq:homogeneous}可知 $z = C\eu^{ax}$，从而得公式、
存在性与唯一性。若 $a < 0$，则 $\eu^{ax} \to 0$，
即当 $x \to +\infty$ 时 $y(x) \to -\frac ba$。
\end{proof}

\begin{method}[由初始条件求解 $y' = ay + b$]\label{met:g12:diffeq:solve}
\begin{enumerate}
  \item 求平衡解 $y_p = -\frac{b}{a}$（即解 $y' = 0$）。
  \item 写出通解 $y = C\eu^{ax} + y_p$。
  \item 由初始条件确定 $C$。
  \item 结合物理直观检验长期性态（解是否
        收敛于平衡解？）。
\end{enumerate}
同样的策略 --- \emph{特解 + 齐次方程
的通解} --- 可推广至 $y' = ay + f(x)$，
见\cref{prop:g12:diffeq:general}。
\end{method}

\begin{example}[牛顿冷却定律]\label{ex:g12:diffeq:cooling}
一杯温度为 $80\,^\circ$C 的咖啡置于温度为 $20\,^\circ$C 的房间中。牛顿冷却定律指出
温度 $T$ 满足 $T' = -k(T - 20)$，其中 $k > 0$ 为某常数，
\emph{即} $T' = -kT + 20k$。平衡解为 $20$，且
$T(t) = 20 + 60\,\eu^{-kt}$：咖啡按指数规律快速冷却至
室温。
\end{example}

\begin{omfigure}
\begin{tikzpicture}
\begin{axis}[omaxis,
  width=9.2cm, height=5.8cm,
  xlabel={$t$}, ylabel={$T$},
  xmin=0, xmax=10.5, ymin=0, ymax=88,
  xtick={5,10}, ytick={20,50,80}]
  \addplot[omProp, dashed, domain=0:10.4] {20};
  \addplot[omDef, thick, domain=0:10.2, samples=60] {20 + 60*exp(-0.5*x)};
  \addplot[omDef!50, thick, domain=0:10.2, samples=60] {20 + 30*exp(-0.5*x)};
  \addplot[omDef!50, thick, domain=0:10.2, samples=60] {20 - 15*exp(-0.5*x)};
  \node[omProp, font=\small, anchor=south west] at (axis cs:6.6,20.8)
    {平衡解};
\end{axis}
\end{tikzpicture}

{\small 无论初始温度如何，$T' = -k(T - 20)$ 的所有解
都指数收敛于平衡解 $T = 20$。}
\end{omfigure}

\section{方程 \texorpdfstring{$y' = ay + f(x)$}{y' = ay + f(x)}}

\begin{proposition}[解集的结构]\label{prop:g12:diffeq:general}
设 $f$ 在区间 $I$ 上连续，$y_p$ 为
下述方程的一个特解
\[
y' = ay + f(x)
\]
在 $I$ 上。则 $I$ 上的解恰为
函数 $y = C\eu^{ax} + y_p$，$C \in \R$。
\end{proposition}

\begin{proof}
同\cref{thm:g12:diffeq:affine}，$y$ 为解当且仅当
$z = y - y_p$ 满足
$z' = (ay + f) - (ay_p + f) = az$，当且仅当 $z = C\eu^{ax}$。
\end{proof}

\begin{method}[猜测特解]\label{met:g12:diffeq:particular}
求一个与 $f$ \emph{形式相同}的特解：
\begin{itemize}
  \item 若 $f$ 为 $n$ 次多项式，则试取 $n$ 次多项式；
  \item 若 $f(x) = \alpha\,\eu^{kx}$ 且 $k \neq a$，则试取 $y_p = \beta\,\eu^{kx}$；
  \item 若 $f(x) = \alpha\,\eu^{ax}$（共振情形），则试取
        $y_p = \beta x\,\eu^{ax}$。
\end{itemize}
代入方程，并比较系数以确定待定常数。
\end{method}

\begin{example}\label{ex:g12:diffeq:particular}
求解 $y' = 2y + 4x$。试取 $y_p = \alpha x + \beta$：
$\alpha = 2(\alpha x + \beta) + 4x$ 对所有 $x$ 成立，可得
$2\alpha + 4 = 0$ 且 $\alpha = 2\beta$，故 $\alpha = -2$，$\beta = -1$。
通解为 $y = C\eu^{2x} - 2x - 1$。
\end{example}

\section{练习}

\begin{exercise}[$\star$]\label{exo:g12:diffeq:1}
在 $\R$ 上求解：
(a) $y' = 3y$，且 $y(0) = 2$；\quad
(b) $2y' + y = 0$，且 $y(0) = -1$；\quad
(c) $y' = -y + 5$，且 $y(0) = 0$。
\end{exercise}

\begin{exercise}[$\star$]\label{exo:g12:diffeq:2}
某细菌种群的增长速率与其数量成正比，
每 $3$ 小时翻一番。写出该种群数量
$N(t)$ 所满足的微分方程，并确定比例系数。
\end{exercise}

\begin{exercise}[$\star$]\label{exo:g12:diffeq:3}
验证 $y_p(x) = x\,\eu^{x}$ 是方程 $y' = y + \eu^x$ 的解，
并求其在 $\R$ 上的所有解。
\end{exercise}

\begin{exercise}[$\star\star$]\label{exo:g12:diffeq:4}
求解 $y' = -2y + \eu^{x}$，满足 $y(0) = 1$。
（提示：求形如 $\beta\,\eu^{x}$ 的特解。）
\end{exercise}

\begin{exercise}[$\star\star$]\label{exo:g12:diffeq:5}
碳-14 的半衰期为 $5730$ 年。一件考古样本
中的碳-14 含量为活体生物的 $60\%$。估算该样本的年龄。
\end{exercise}

\begin{exercise}[$\star\star$]\label{exo:g12:diffeq:6}
水箱中盛有 $100$ L 纯水。每升含盐 $0.2$ kg 的盐水
以 $5$ L/min 的速率流入，混合液（保持均匀）以相同速率流出。
设 $m(t)$ 为时刻 $t$（以分钟计）时水箱中盐的质量。
\begin{enumerate}
  \item 说明 $m' = 1 - \dfrac{m}{20}$ 成立。
  \item 求解，并求 $m(t)$ 当 $t \to +\infty$ 时的极限，
        并解释其实际意义。
\end{enumerate}
\end{exercise}

\begin{exercise}[$\star\star$]\label{exo:g12:diffeq:7}
质量为 $80$ kg 的跳伞运动员受重力作用下落，
（$g = 9.8\ \text{m/s}^2$），空气阻力与速率成正比，使得
其速率满足 $v' = g - \frac{k}{m}v$，其中 $k = 16$ kg/s。
\begin{enumerate}
  \item 在 $v(0) = 0$ 的条件下求解方程。
  \item 计算\emph{终端速度} $\lim_{t\to+\infty} v(t)$，以及
        达到其 $95\%$ 所需的时间。
\end{enumerate}
\end{exercise}

\begin{exercise}[$\star\star\star$]\label{exo:g12:diffeq:8}
\emph{（逻辑斯谛方程）}一个种群 $y(t) \in \intoo{0}{1}$（以最大种群数量的分数表示）
满足下列方程：
\[
y' = y(1 - y).
\]
\begin{enumerate}
  \item 令 $z = \dfrac{1}{y}$，证明 $z$ 满足线性方程
        $z' = -z + 1$。
  \item 求 $z$，再求 $y$，其中 $y(0) = \frac{1}{10}$。
  \item 证明 $y(t) \to 1$（当 $t \to +\infty$ 时），
        并大致画出解曲线的形状。
\end{enumerate}
\end{exercise}

\begin{exercise}[$\star\star\star$]\label{exo:g12:diffeq:9}
设 $y$ 为 $y' = ay + b$ 的一个解，令 $u_n = y(n)$，其中 $n \in \N$。
证明 $(u_n)$ 满足一阶线性递推关系
$u_{n+1} = q u_n + r$，并把 $q$ 和 $r$ 用 $a$ 和 $b$ 表示。
条件 $\abs{q} < 1$ 对该微分方程
意味着什么？
\end{exercise}

\section{问题：事物内部的时钟}

\begin{problem}[{周末问题 --- 碳-14 为洞穴定年，
牛顿冷却定律与一桩命案，一个小方程穿上
四套戏服}]
\label{pb:g12:diffeq:1}
微分方程是变化的规律；求解它，就把这条规律
变成一只\emph{钟}。同一个小方程
$y' = ay + b$（\cref{thm:g12:diffeq:affine}）在
史前木炭、变凉的咖啡、下落的雨滴和
医院的输液中滴答走动 --- 读取这些钟，正是本题的任务：
它要为拉斯科洞穴壁画定年，确定死亡时间，
并像一位好科学家那样检验自身的假设。

\medskip
\noindent\textbf{第一部分 --- 熟练运算。}
\begin{enumerate}
  \item 求解 $y' = 3y$，$y(0) = 2$；再求
        $y' = -2y + 6$，$y(0) = 0$ 的解
        （\cref{met:g12:diffeq:solve}）。
  \item 重新推导
        \cref{thm:g12:diffeq:homogeneous} 背后的唯一性技巧：若 $y' = ay$，计算
        $y(t)\,\eu^{-at}$ 的导数，并得出
        \emph{每个}解都是 $C\eu^{at}$。
  \item 对 $y' = -2y + 6$，求平衡解，并
        描述其余每个解在
        $t \to +\infty$ 时的走向。（不动点来自
        \cref{pb:g11:seq:1} 的递推关系，只是改为连续情形。）
  \item 对衰变规律 $y' = -\frac{y}{\tau}$，证明
        半衰期为 $t_{1/2} = \tau \ln 2$。
  \item 画出（或描述）方程
        $y' = -2y + 6$ 的解族：从高于 $3$ 处出发的解会怎样？
        从低于 $3$ 处出发呢？恰为 $3$ 呢？
\end{enumerate}

\medskip
\noindent\textbf{第二部分 --- 碳-14。}
活体组织中放射性碳-14 的比例保持恒定；
死亡后摄入停止，存量开始衰减：
$N' = -\lambda N$，半衰期为 $5\,730$ 年。
\begin{enumerate}[resume]
  \item 计算 $\lambda$（单位：每年）。
  \item 某骨头保留了活体比例的 $20\,\%$，
        它的年龄是多少？
  \item 取自拉斯科洞穴壁画的木炭保留了
        约 $15\,\%$，请为这些壁画定年。
  \item 在阿尔卑斯冰川中发现的冰人奥茨，测得
        约 $53\,\%$，请为他定年（考古学家认为约
        $5\,300$ 年 --- 你算得怎样？）。
  \item 为何碳-14 无法测定恐龙的年代？计算
        经过十个半衰期后的剩余比例，
        再写出 $65$ 百万年后的剩余比例（以 $2$ 的幂表示），
        并作出结论。
  \item 误差范围：若奥茨的 $53\,\%$ 仅精确到
        $\pm 1\,\%$，求年龄的取值范围。实验精度
        能换来怎样的定年精度？
\end{enumerate}

\medskip
\noindent\textbf{第三部分 --- 牛顿冷却与一桩命案。}
温度为 $T$ 的物体处于温度恒为 $T_a$ 的环境中，
其按 $T' = -k\,(T - T_a)$ 的规律降温。
\begin{enumerate}[resume]
  \item 求解该方程（令 $z = T - T_a$）：
        \[
        T(t) = T_a + (T_0 - T_a)\,\eu^{-kt} .
        \]
  \item 咖啡以 $90\,^\circ$C 倒出，置于 $20\,^\circ$C 的房间中，
        读数为 $70\,^\circ$C，此时已过去 $5$ 分钟。求 $k$，
        再求降到可饮用的
        $55\,^\circ$C 所需的等待时间。
  \item 牛奶问题：要想十分钟后喝到尽可能热的咖啡，
        冷牛奶应现在加入，还是
        等到最后时刻才加入？请用方程作答（
        加入牛奶会对温差 $T - T_a$ 产生什么影响，
        温差又如何驱动热量的散失？）。
  \item 法医问题：一具尸体于午夜被发现，
        温度为 $30\,^\circ$C，身处 $20\,^\circ$C 的房间中；一小时后
        读数为 $28\,^\circ$C。假设死亡时温度为 $37\,^\circ$C，
        由两次测量求出 $k$，
        再求死亡时间。
  \item 检验验尸官的时钟：指出三条与现实不符的模型假设，
        并说明每一条假设被违背时会使死亡时间的估计
        偏早还是偏晚。
\end{enumerate}

\medskip
\noindent\textbf{第四部分 --- 四套戏服与一条 S 形曲线。}
\begin{enumerate}[resume]
  \item \cref{exo:g12:diffeq:8} 中的逻辑斯谛方程，
        $y' = y(1 - y)$，$y(0) = 0.1$：利用该练习的
        代换，推导出
        $y(t) = \dfrac{1}{1 + 9\eu^{-t}}$，并求
        拐点 $y = \frac12$ 对应的时刻 --- 这正是
        \cref{pb:g12:deriv:1} 中流行病学家所说的那个日期，
        现在可以算出来了。
  \item 雨滴满足 $v' = 10 - \frac v2$（重力减去阻力），
        $v(0) = 0$。求终端速度，解出
        $v(t)$，并计算雨滴何时达到
        终端速度的 $95\,\%$。
  \item 医院的输液以恒定速率输入药物，
        而人体按比例将药物排出：
        $c' = 4 - 0.5c$，$c(0) = 0$。求稳态浓度，
        以及达到稳态浓度一半所需的时间。用一句话
        指出它与
        \cref{exo:g12:diffeq:9} 中贷款递推关系的相似之处。
  \item 压轴题 --- 一个方程，四套戏服：衰变、冷却、
        下落、给药，都是 $y' = ay + b$，只是符号各不相同，
        名称也不同。背出本题重复了四次的建模循环
        （规律 $\to$ 方程 $\to$ 解 $\to$
        标定 $\to$ 预测 $\to$ 检验），并指出
        需要 $y''$ 的现象 --- 以及哪一道题
        已经遇到过它。
\end{enumerate}
\end{problem}
